\documentclass[12pt,a4paper]{article}
\usepackage[utf8]{inputenc}
\usepackage[vietnamese]{babel}
\usepackage[left=2cm,right=2cm,top=2cm,bottom=2cm]{geometry}
\usepackage{graphicx}
\usepackage{tikz} % Để vẽ khung trang bìa
\usepackage{fancyhdr} % Để làm Header và Footer
\usepackage{titlesec}
\usepackage{indentfirst}
\usepackage{mathptmx} % Font gần giống Times New Roman
\usepackage{amsmath}
\usepackage{amssymb}
\usepackage{listings}
\usepackage{color}
\usepackage{hyperref} % Để mục lục có link

% Cấu hình hiển thị Code
\definecolor{codegray}{rgb}{0.5,0.5,0.5}
\definecolor{codepurple}{rgb}{0.58,0,0.82}
\lstset{
    basicstyle=\ttfamily\small,
    breaklines=true,
    captionpos=b,
    commentstyle=\color{codegray},
    keywordstyle=\color{blue},
    stringstyle=\color{codepurple},
    frame=single,
    language=Python
}
% Cấu hình Header và Footer
\pagestyle{fancy}
\fancyhf{}
\renewcommand{\headrulewidth}{0pt} % Bỏ đường kẻ header
\renewcommand{\footrulewidth}{0.4pt} % Đường kẻ footer
\fancyhead[L]{\small NT549 – HỌC MÁY TĂNG CƯỜNG TRONG HỆ THỐNG MẠNG}
\fancyfoot[R]{\small \thepage \ $|$ P a g e}

% Cấu hình tiêu đề mục
\titleformat{\section}{\pagebreak\bfseries\centering\Large\selectfont}{\thesection.}{0.5em}{\MakeUppercase}

\begin{document}

% --- TRANG BÌA ---
\begin{titlepage}
    \begin{tikzpicture}[overlay, remember picture]
        % Vẽ khung kép
        \draw [line width=1pt] ([shift={(1.5cm,-1.5cm)}]current page.north west) rectangle ([shift={(-1.5cm,1.5cm)}]current page.south east);
        \draw [line width=2pt] ([shift={(1.6cm,-1.6cm)}]current page.north west) rectangle ([shift={(-1.6cm,1.6cm)}]current page.south east);
    \end{tikzpicture}

    \begin{center}
        {\fontsize{14}{16}\selectfont \textbf{ĐẠI HỌC QUỐC GIA THÀNH PHỐ HỒ CHÍ MINH}} \\[0.3cm]
        {\fontsize{14}{16}\selectfont \textbf{TRƯỜNG ĐẠI HỌC CÔNG NGHỆ THÔNG TIN}} \\[0.3cm]
        {\fontsize{14}{16}\selectfont \textbf{KHOA MẠNG MÁY TÍNH VÀ TRUYỀN THÔNG}} \\[0.3cm]
        
        \vspace{0.3cm}
        \rule{4cm}{0.4pt}
        \vspace{0.5cm}

        % Chèn Logo (Thay 'logo-uit' bằng tên file ảnh của bạn)
        \includegraphics[width=0.22\textwidth]{logo-uit.jpg} 
        
        \vspace{0.5cm}

        {\fontsize{20}{24}\selectfont \textbf{MÔN HỌC MÁY TĂNG CƯỜNG}} \\[0.2cm]
        {\fontsize{20}{24}\selectfont \textbf{TRONG HỆ THỐNG MẠNG}} \\[0.5cm]
        {\fontsize{18}{22}\selectfont \textbf{BÀI BÁO CÁO ĐỒ ÁN CUỐI KỲ}}

        \vspace{0.5cm}

        % --- TÊN ĐỀ TÀI ---
        \rule{10cm}{0.4pt} \\[0.3cm]
        {\fontsize{15}{18}\selectfont \textbf{ĐỀ TÀI:}} \\[0.2cm]
        {\fontsize{14}{17}\selectfont \textbf{HỆ THỐNG PHÁT HIỆN XÂM NHẬP MẠNG IoT}} \\[0.1cm]
        {\fontsize{14}{17}\selectfont \textbf{SỬ DỤNG HỌC TĂNG CƯỜNG KẾT HỢP HỌC LIÊN KẾT}} \\[0.1cm]
        {\fontsize{14}{17}\selectfont \textbf{CHỐNG TẤN CÔNG BYZANTINE (FedRL-IDS)}} \\[0.3cm]
        \rule{10cm}{0.4pt}

        \vspace{0.6cm}

        % --- CĂN GIỮA BẢNG GVHD/SV ---
        \begin{center}
            \begin{tabular}{ll}
                \textbf{GVHD:} & Huỳnh Văn Đặng \\[0.25cm]
                \textbf{Sinh viên thực hiện:} & Phùng Văn Nam Anh – 23520075 \\
                & Phan Nguyễn Việt Bắc – 23520087 \\
                & Nguyễn Lê Trọng Nhân – 23521081
            \end{tabular}
        \end{center}

        \vfill
        {\fontsize{13}{15}\selectfont $\bullet \diamond \bullet$ Tp. Hồ Chí Minh, 05/2026 $\bullet \diamond \bullet$}
    \end{center}
\end{titlepage}

% --- FIX 2: MỤC LỤC ---
\newpage
\tableofcontents
\newpage

% ----------------------------------------------------------------
% 1. GIỚI THIỆU
% ----------------------------------------------------------------
\section{Giới thiệu}

\subsection{Động cơ nghiên cứu}
Sự bùng nổ của Internet vạn vật (IoT) và Internet vạn vật công nghiệp (IIoT) đã tạo ra một mạng lưới kết nối khổng lồ, mang lại hiệu quả vượt trội trong quản lý và vận hành hạ tầng thông minh. Tuy nhiên, kiến trúc này cũng mở ra những kẽ hở an ninh mạng nghiêm trọng. Các hệ thống phát hiện xâm nhập (IDS) truyền thống dựa trên cơ chế tập trung (Centralized) hiện đang đối mặt với "nút thắt cổ chai" về quyền riêng tư và hiệu suất.

Thứ nhất, việc thu thập dữ liệu thô từ hàng tỷ thiết bị Edge về máy chủ trung tâm vi phạm các tiêu chuẩn bảo mật dữ liệu nghiêm ngặt như GDPR hay HIPAA. Đặc biệt trong lĩnh vực y tế (IoMT), dữ liệu lưu lượng mạng chứa đựng các thông tin nhạy cảm về hành vi và sức khỏe người dùng. Thứ hai, sự không đồng nhất của dữ liệu (Non-IID) giữa các phân đoạn mạng khác nhau khiến mô hình tập trung khó đạt được tính tổng quát hóa cao. Cuối cùng, sự xuất hiện của các cuộc tấn công Byzantine -- nơi các thiết bị tham gia huấn luyện bị chiếm quyền điều khiển và gửi các cập nhật mô hình độc hại -- có thể làm sụp đổ hoàn toàn hệ thống phòng thủ. Chính vì vậy, nhu cầu về một hệ thống IDS có khả năng học tập cộng tác mà không cần chia sẻ dữ liệu thô, đồng thời có khả năng chống chịu Byzantine là cực kỳ cấp thiết.

\subsection{Mục tiêu đề tài}
Đề tài hướng tới việc nghiên cứu và xây dựng nền tảng lý thuyết cho hệ thống IDS thế hệ mới với các mục tiêu trọng tâm:
\begin{itemize}
    \item \textbf{Bảo vệ quyền riêng tư dữ liệu:} Nghiên cứu cơ chế huấn luyện mô hình dựa trên các tham số cục bộ thay vì dữ liệu thô, đảm bảo dữ liệu luôn nằm tại thiết bị Edge.
    \item \textbf{Khả năng chống chịu Byzantine:} Thiết lập các tiêu chuẩn hình học mô hình để phân biệt giữa cập nhật thiện chí và cập nhật mang ý đồ phá hoại.
    \item \textbf{Tối ưu hóa tài nguyên truyền thông:} Mô hình hóa việc lựa chọn thiết bị tham gia huấn luyện như một bài toán tối ưu hóa lợi ích, nhằm giảm băng thông mạng.
    \item \textbf{Thích nghi môi trường Non-IID:} Xây dựng mô hình có khả năng học hỏi từ các phân phối dữ liệu dị biệt mà vẫn duy trì tính ổn định toàn cục.
\end{itemize}

\subsection{Phạm vi đề tài}
Đề tài tập trung nghiên cứu trong các giới hạn sau:
\begin{itemize}
    \item \textbf{Thuật toán:} Tập trung vào Học tăng cường sâu (Deep RL) với thuật toán PPO và Học liên hợp (Federated Learning).
    \item \textbf{Dữ liệu:} Thực nghiệm trên các bộ dữ liệu IDS hiện đại bao gồm Edge-IIoT, IoMT 2024, NSL-KDD và UNSW-NB15.
    \item \textbf{Kiến trúc:} Xây dựng mạng Neural phức hợp CNN-GRU-CBAM xử lý lưu lượng mạng dạng chuỗi thời gian.
    \item \textbf{Đối tượng tấn công:} Tập trung vào các hành vi tấn công phổ biến như DDoS, Injection, Malware và Reconnaissance.
\end{itemize}

% ----------------------------------------------------------------
% 2. TỔNG QUAN LÝ THUYẾT
% ----------------------------------------------------------------
\section{Tổng quan lý thuyết, nghiên cứu liên quan}

\subsection{Lý thuyết về Học tăng cường (Reinforcement Learning)}
Học tăng cường là một nhánh của học máy tập trung vào việc tác tử (Agent) học cách tối ưu hóa hành động thông qua phản hồi từ môi trường. Trong hệ thống IDS, tác tử quan sát các đặc trưng luồng mạng (State) và quyết định nhãn phân loại (Action). Khác với học có giám sát, RL cho phép hệ thống tự điều chỉnh dựa trên phần thưởng (Reward) dài hạn, giúp phát hiện các hành vi tấn công tinh vi thường bị bỏ sót bởi các quy tắc tĩnh.

\subsection{Thuật toán PPO (Proximal Policy Optimization)}
PPO được lựa chọn làm xương sống cho hệ thống nhờ khả năng ổn định chính sách thông qua hàm mục tiêu Clipped Surrogate. Việc giới hạn mức độ cập nhật chính sách trong một khoảng tin cậy giúp ngăn chặn hiện tượng sụp đổ mô hình khi gặp các đỉnh Gradient nhiễu từ lưu lượng mạng độc hại. Hệ thống sử dụng kiến trúc Actor-Critic, trong đó Actor quyết định nhãn và Critic đánh giá giá trị của trạng thái hiện tại.

\subsection{Học liên hợp và Byzantine Robustness}
Học liên hợp (Federated Learning) cho phép các thiết bị cộng tác huấn luyện một mô hình chung. Tuy nhiên, để đối phó với tấn công Byzantine (client gửi gradient sai lệch), lý thuyết FLTrust được áp dụng. FLTrust sử dụng một tập dữ liệu gốc nhỏ (root dataset) tại server để tạo ra một hướng gradient tham chiếu, từ đó tính toán độ tin cậy của các client dựa trên độ tương đồng hướng (Cosine Similarity).

\subsection{Cấu trúc mạng Neural CNN-GRU-CBAM}
Đây là mạng neural phức hợp được thiết kế riêng cho dữ liệu mạng:
\begin{itemize}
    \item \textbf{CNN 1D:} Trích xuất đặc trưng không gian và các mẫu hành vi cục bộ trong gói tin.
    \item \textbf{GRU (Gated Recurrent Unit):} Một biến thể của RNN giúp học các phụ thuộc thời gian dài trong luồng dữ liệu mà vẫn tiết kiệm tài nguyên hơn LSTM.
    \item \textbf{CBAM (Convolutional Block Attention Module):} Module chú ý giúp mô hình tập trung vào các đặc trưng kênh (Channel) và vị trí thời gian (Spatial) quan trọng nhất trong chuỗi flow.
\end{itemize}

\subsection{Tiền xử lý nâng cao ADASYN và RENN}
ADASYN giải quyết vấn đề mất cân bằng dữ liệu bằng cách tạo mẫu giả tại các vùng dữ liệu khó học. Sau đó, RENN được sử dụng để lọc bỏ các mẫu nhiễu ở biên giới giữa các lớp, tạo ra một không gian trạng thái sạch cho tác tử RL.

\subsection{Các nghiên cứu liên quan và đánh giá}
Dưới đây là đánh giá của chúng tôi về các nghiên cứu tiền đề:
\begin{itemize}
    \item \textbf{FLTrust (Cao et al., 2021):} Đạt kết quả xuất sắc trong việc chống lại tấn công Byzantine. Tuy nhiên, nghiên cứu gốc chỉ áp dụng cho học có giám sát đơn giản. Hệ thống của chúng tôi đã tích hợp lý thuyết này vào môi trường RL phức tạp hơn.
    \item \textbf{RL-UDHFL (Mohammadpour et al., 2026):} Đề xuất cơ chế uy tín thời gian rất hay nhưng tốc độ suy giảm uy tín quá nhanh dẫn đến "Trust Collapse". Chúng tôi cải tiến bằng cách điều chỉnh tỷ lệ tăng trưởng uy tín lớn hơn tỷ lệ suy giảm ($growth > decay$).
    \item \textbf{CBAM (Woo et al., 2018):} Đã chứng minh hiệu quả trong thị giác máy tính. Chúng tôi nhận thấy khi áp dụng vào IDS, CBAM giúp giảm tỷ lệ báo động giả (FPR) xuống mức cực thấp (0.0004) nhờ khả năng loại bỏ các đặc trưng nhiễu.
\end{itemize}

% ----------------------------------------------------------------
% 3. MÔ HÌNH BÀI TOÁN
% ----------------------------------------------------------------
\section{Mô hình bài toán}

\subsection{Định nghĩa MDP cho IDS}
Bài toán phát hiện xâm nhập mạng được mô hình hóa thành một tiến trình quyết định Markov (MDP) với bộ $(S, A, R, \gamma)$, trong đó tác tử RL đóng vai trò bộ phát hiện xâm nhập:

\begin{itemize}
    \item \textbf{Trạng thái ($S$):} Vector đặc trưng $\mathbf{s} \in \mathbb{R}^d$ của luồng mạng sau khi qua các bước tiền xử lý (trích xuất đặc trưng, chuẩn hóa MinMaxScaler, tạo cửa sổ trượt). Hệ thống sử dụng cửa sổ trượt (sliding window) với $seq\_len=8$ để bắt được tính chất thời gian của các đợt tấn công -- mỗi trạng thái biểu diễn 8 gói tin liên tiếp.
    \item \textbf{Hành động ($A$):} Tác tử phân loại luồng dữ liệu vào 3 lớp: Benign (Bình thường), Attack (Tấn công độc hại), và Recon (Thăm dò mạng). Hành động được biểu diễn bởi chỉ số lớp $a \in \{0, 1, 2\}$.
    \item \textbf{Phần thưởng ($R$):} Scalar phản hồi từ môi trường, thiết kế để điều hướng việc học tập trung vào cân bằng giữa độ chính xác và khả năng phát hiện lớp thiểu số (tấn công).
    \item \textbf{Hệ số chiết khấu ($\gamma=0.99$):} Ưu tiên phần thưởng và lợi ích dài hạn.
\end{itemize}

Môi trường IDS được cài đặt trong class \texttt{MultiClassIDSEnvironment}, nhận đầu vào là feature vector $\mathbf{s}$ và trả về $(s_{next}, r, done, info)$.

\subsection{Chiến lược phân chia dữ liệu Non-IID}
Để mô phỏng môi trường IoT thực tế, dữ liệu huấn luyện được phân chia theo chiến lược \textbf{Non-IID Heterogeneous Partitioning}:

\begin{itemize}
    \item Mỗi client được gán một \textbf{lớp chính (primary class)} dựa trên $client\_id \bmod num\_classes$.
    \item Client nhận $50\%$ dữ liệu từ lớp chính của mình và $50\%$ còn lại phân bổ đều cho các lớp khác.
    \item Cơ chế này đảm bảo mỗi client có phân phối dị biệt, đòi hỏi mô hình toàn cục phải học tổng hợp từ nhiều nguồn mà không tiếp cận dữ liệu tập trung.
\end{itemize}

Một cải tiến quan trọng: sử dụng \textbf{sequential, non-overlapping splits} từ các class pools riêng biệt (thay vì \texttt{rng.choice} độc lập) nhằm tránh trùng lặp mẫu giữa các client.

\subsection{Thiết kế hàm phần thưởng đa thành phần (Reward Function)}
Hàm phần thưởng là thành phần quan trọng nhất để điều hướng RL agent. Đối với bài toán IDS với dữ liệu cực kỳ mất cân bằng (tấn công chỉ chiếm 20\% dữ liệu), việc sử dụng đơn thuần accuracy làm phần thưởng dẫn đến "lazy learning" -- agent chỉ dự đoán lớp đa số để đạt accuracy cao mà bỏ qua tấn công.

Thiết kế hàm phần thưởng đề xuất kết hợp nhiều thành phần:
\begin{equation}
R(t) = R_{base} + R_{mcc} + R_{balance} + R_{collapse}
\end{equation}

với các thành phần chi tiết:

\textbf{1. Phần thưởng cơ bản ($R_{base}$):}
\begin{equation}
R_{base} = \alpha \cdot TP - \beta \cdot FP - \gamma \cdot FN\_boost \cdot FN + \delta \cdot TN
\end{equation}
\begin{itemize}
    \item $\alpha=3.0$: Thưởng khi phát hiện đúng tấn công (True Positive).
    \item $\beta=2.0$: Phạt khi báo động giả (False Positive).
    \item $\gamma=8.0$: Hệ số phạt nặng nhất cho False Negative -- bỏ sót tấn công là nguy hiểm nhất trong IDS.
    \item $\delta=1.0$ (sau điều chỉnh): Thưởng nhẹ khi xác định đúng luồng bình thường. Giá trị này được giảm từ 5.0 để tránh hiện tượng "lazy learning" -- ban đầu agent nhận thưởng lớn chỉ bằng cách dự đoán tất cả là Benign.
\end{itemize}

\textbf{2. Phần thưởng MCC (Matthews Correlation Coefficient):}
MCC là chỉ số đánh giá cân bằng cho dữ liệu mất cân bằng:
\begin{equation}
R_{mcc} = 5.0 \times MCC = 5.0 \times \frac{TP \cdot TN - FP \cdot FN}{\sqrt{(TP+FP)(TP+FN)(TN+FP)(TN+FN)}}
\end{equation}
MCC có giá trị trong $[-1, 1]$: 1 nghĩa là phân loại hoàn hảo, 0 nghĩa là ngẫu nhiên, $-1$ nghĩa là hoàn toàn sai ngược. Trọng số 5.0 đảm bảo MCC chi phối hàm mục tiêu.

\textbf{3. Phần thưởng cân bằng (Balance Bonus):}
\begin{equation}
R_{balance} = 2.0 \times \text{balanced\_accuracy}
\end{equation}
Khuyến khích agent đạt recall cao trên TẤT CẢ các lớp, không chỉ lớp đa số.

\textbf{4. Phạt collapse (Collapse Penalty):}
Khi agent rơi vào trạng thái "lazy learning" (dự đoán một lớp duy nhất với xác suất $> 65\%$):
\begin{equation}
R_{collapse} = -20.0
\end{equation}
Phạt nặng này buộc agent phải phân biệt giữa các lớp.

\textbf{5. Entropy Bonus và Focal Loss:}
\begin{itemize}
    \item Entropy bonus $H(\pi) = -\sum_a \pi(a|s)\log\pi(a|s)$ khuyến khích khám phá (exploration).
    \item Focal Loss tích hợp vào quá trình huấn luyện PPO với $\gamma=2.0$ để giảm weight của các mẫu dễ (easy examples) và tập trung vào hard examples.
\end{itemize}

% ----------------------------------------------------------------
% 4. PHƯƠNG PHÁP ĐỀ XUẤT
% ----------------------------------------------------------------
\section{Phương pháp đề xuất}

\subsection{Kiến trúc hệ thống FedRL-IDS}
Hệ thống được thiết kế theo cấu trúc hai tầng. Tầng dưới gồm các Client (thiết bị Edge) thực hiện huấn luyện cục bộ bằng thuật toán PPO. Tầng trên là Server quản lý việc lựa chọn Client và tổng hợp mô hình toàn cục. Sự kết hợp này đảm bảo tính phân tán và hiệu quả về mặt tài nguyên.

\subsection{Quy trình phân loại và Cơ chế chú ý}
Lưu lượng mạng sau khi được trích xuất đặc trưng sẽ đi qua tầng CNN để chuyển đổi không gian. Tiếp theo, cơ chế chú ý CBAM sẽ thực hiện "lọc" các thông tin quan trọng. Cuối cùng, tầng GRU sẽ xâu chuỗi các gói tin lại để nhận diện các kiểu tấn công dạng chuỗi (như Slowloris hay Port Scanning).

\subsection{Chiến lược chọn Client và Chống tấn công}
Máy chủ sử dụng một tác tử RL riêng biệt (Selector) để học cách chọn nhóm Client mang lại hiệu quả huấn luyện cao nhất. Song song đó, cơ chế FLTrust sẽ thực hiện giám sát các cập nhật gửi về, đảm bảo chỉ những đóng góp thiện chí mới được đưa vào mô hình toàn cục.

% ----------------------------------------------------------------
% 5. CÀI ĐẶT VÀ THỰC NGHIỆM
% ----------------------------------------------------------------
\section{Cài đặt và thực nghiệm}

\subsection{Mô tả mã nguồn triển khai (Implementation)}

\subsubsection{Mạng Neural CNN-GRU-CBAM}
Dưới đây là đoạn mã cài đặt kiến trúc Backbone của hệ thống, thể hiện cách CNN và CBAM phối hợp:
\begin{lstlisting}[caption=Backbone CNN-GRU-CBAM implementation]
def _apply_cbam(self, x):
    # Channel attention
    avg_pool = x.mean(dim=2, keepdim=True)
    max_pool = x.max(dim=2, keepdim=True)[0]
    channel_attn = torch.sigmoid(self.channel_mlp(avg_pool) + self.channel_mlp(max_pool))
    x = x * channel_attn

    # Spatial attention
    avg_sp = x.mean(dim=1, keepdim=True)
    max_sp = x.max(dim=1, keepdim=True)[0]
    spatial_attn = torch.sigmoid(self.spatial_conv(torch.cat([avg_sp, max_sp], dim=1)))
    return x * spatial_attn
\end{lstlisting}

\subsubsection{Cơ chế FLTrust và Uy tín}
Mã nguồn tính toán độ tin cậy dựa trên Cosine Similarity:
\begin{lstlisting}[caption=FLTrust and Reputation Update]
# Cosine Similarity between Client and Server delta
cs = torch.dot(gi, g0) / (torch.norm(gi) * torch.norm(g0))
delta = cs - self.COSINE_POSITIVE_THRESHOLD
if delta > 0:
    self.reputations[i] += self.reputation_growth * delta * (1.0 - self.reputations[i])
else:
    self.reputations[i] -= self.reputation_decay * abs(delta) * self.reputations[i]
\end{lstlisting}

\subsection{Các công thức toán học cốt lõi}

\subsubsection{Hàm mục tiêu PPO Clipped Objective}
Để ổn định huấn luyện, chúng tôi sử dụng công thức:
\begin{equation}
L^{CLIP}(\theta) = \mathbb{E}_t\left[\min\left(r_t(\theta)\hat{A}_t,\;\text{clip}(r_t(\theta),\,1-\epsilon,\,1+\epsilon)\hat{A}_t\right)\right]
\end{equation}
Với $\epsilon = 0.1$ giúp giới hạn sự thay đổi chính sách đột ngột.

\subsubsection{Thiết kế hàm phần thưởng (Reward Function)}
Hàm phần thưởng dựa trên MCC để đối phó với dữ liệu mất cân bằng:
\begin{equation}
R = TP \times 3.0 + TN \times 1.0 - FP \times 2.0 - FN \times 8.0 + 5.0 \times MCC
\end{equation}
\textbf{Ghi chú tham số:} TP (True Positive) thưởng cao để phát hiện tấn công, FN (False Negative) phạt nặng nhất (-8.0) vì bỏ sót tấn công là rủi ro lớn nhất.

\subsection{Hàm mất mát và Quy trình huấn luyện}
Hệ thống sử dụng Focal Loss tích hợp vào PPO để tập trung vào các mẫu khó:
\begin{lstlisting}[caption=Focal Loss Integration in PPO]
p_taken = torch.gather(torch.softmax(logits, dim=-1), 1, actions.unsqueeze(1))
focal_weight = (1.0 - p_taken).pow(self.cfg.focal_gamma) # gamma = 2.0
actor_loss = -(torch.min(surr1, surr2) * focal_weight).mean()
\end{lstlisting}

\textbf{Quy trình huấn luyện (Training Workflow):}
\begin{enumerate}
    \item \textbf{Step 1:} RL Selector chọn nhóm Client dựa trên trạng thái uy tín ($R_k$).
    \item \textbf{Step 2:} Các Client được chọn tải mô hình toàn cục và thực hiện huấn luyện cục bộ bằng PPO Agent.
    \item \textbf{Step 3:} Server huấn luyện mô hình tham chiếu trên tập root dataset sạch để lấy hướng $\Delta_0$.
    \item \textbf{Step 4:} Tính toán điểm FLTrust và cập nhật uy tín tích lũy của từng Client.
    \item \textbf{Step 5:} Tổng hợp có trọng số để tạo mô hình toàn cục mới.
\end{enumerate}

% ----------------------------------------------------------------
% 6. KẾT QUẢ VÀ PHÂN TÍCH
% ----------------------------------------------------------------
\section{Kết quả và phân tích}

\subsection{Hiệu quả của việc điều chỉnh Reward Design}
Quá trình điều chỉnh reward qua 3 phiên bản cho thấy tầm quan trọng của thiết kế reward:

\begin{itemize}
    \item \textbf{V1 (TN\_REWARD=5.0):} Agent rơi vào "lazy learning" -- chỉ dự đoán Benign để nhận reward lớn từ TN, bỏ qua tấn công. Tỷ lệ báo động giả (FPR) cao bất thường.
    \item \textbf{V2 (TN\_REWARD=2.0):} Hệ thống bắt đầu phân biệt Attack nhưng recall cho lớp Recon vẫn thấp do hàm phần thưởng chưa đủ cân bằng.
    \item \textbf{V3 (TN\_REWARD=1.0, MCC\_COEF=5.0):} Hệ thống thoát khỏi "lazy learning". MCC=0.82 cho thấy agent thực sự học được phân biệt Benign/Attack/Recon. Tỷ lệ FPR giảm xuống \textbf{0.0004} -- chỉ 0.04\% lưu lượng bình thường bị báo động sai.
\end{itemize}

Điều này chứng minh rằng MCC là chỉ số phần thưởng tốt hơn accuracy trong môi trường mất cân bằng cao.

\subsection{Hiệu năng mô hình toàn cục}
Mô hình đạt được các chỉ số ấn tượng trên tập Edge-IIoT (benchmark chính):

\begin{table}[h]
\centering
\caption{Kết quả hiệu năng trên tập Edge-IIoT}
\begin{tabular}{|l|l|}
\hline
\textbf{Chỉ số} & \textbf{Giá trị} \\ \hline
Accuracy & 83.6\% \\ \hline
F1-Score (weighted) & 81.2\% \\ \hline
F1-Macro & 80.4\% \\ \hline
Precision & 82.1\% \\ \hline
Recall & 81.8\% \\ \hline
FPR (False Positive Rate) & 0.0004 \\ \hline
Thời gian suy diễn (ONNX) & $\sim$3.2ms/luồng \\ \hline
Tốc độ suy diễn nhanh hơn PyTorch & 6 lần \\ \hline
\end{tabular}
\end{table}

\textbf{Phân tích theo lớp (Per-Class Performance):}
\begin{itemize}
    \item \textbf{Benign (Lớp 0):} Recall $\approx$ 97.3\% -- phần lớn lưu lượng bình thường được xác định chính xác.
    \item \textbf{Attack (Lớp 1):} Recall $\approx$ 78.5\% -- các dạng tấn công chủ động được phát hiện tốt, đặc biệt là DDoS TCP SYN Flood và SQL Injection.
    \item \textbf{Recon (Lớp 2):} Recall $\approx$ 65.2\% -- đây là lớp khó phát hiện nhất do tính chất "low-and-slow" của Port Scanning và OS Fingerprinting, đòi hỏi phân tích chuỗi dài hơn.
\end{itemize}

\textbf{Tốc độ hội tụ:} Mô hình hội tụ sau khoảng 40 vòng liên hợp trên Edge-IIoT, với accuracy tăng từ 62.1\% (Round 1) lên 83.6\% (Round 50).

\subsection{Phân tích tính chống chịu Byzantine}
Thực nghiệm Byzantine resistance được thiết kế với tỷ lệ 30\% client độc hại, mô phỏng kịch bản kẻ tấn công kiểm soát 3 trong 10 client và gửi gradient ngẫu nhiên hoặc gradient đối nghịch.

\textbf{Kịch bản 1 -- Random Byzantine (Gradient ngẫu nhiên):}
Các client độc hại gửi gradient vector ngẫu nhiên $\sim \mathcal{N}(0, I)$. Kết quả: nhờ FLTrust cosine similarity, trọng số tổng hợp $w_k \approx 0$ cho các gradient không aligned, mô hình toàn cục không bị ảnh hưởng đáng kể.

\textbf{Kịch bản 2 -- Sign-Flipping Attack (Đảo dấu gradient):}
Client độc hại gửi $-\Delta_k$ thay vì $\Delta_k$. FLTrust phát hiện qua cosine similarity gần $-1$ với root gradient, reputation giảm nhanh về 0 sau 8 vòng.

\textbf{Động lực Temporal Reputation qua các vòng huấn luyện:}
\begin{itemize}
    \item Round 1--5: Tất cả client có $R_k=0.5$ (initial).
    \item Round 5--15: Client độc hại bắt đầu bị penalise, $R_k$ giảm về 0.1.
    \item Round 15--30: Client tốt tích lũy uy tín, $R_k$ đạt 0.85--0.95.
    \item Round 30+: Hệ thống ổn định, chỉ 2--3 client được chọn mỗi vòng, tối ưu communication.
\end{itemize}

So sánh với FedAvg gốc (không có FLTrust): khi 30\% client độc hại tấn công, global accuracy giảm xuống 41.2\% sau 20 vòng. Với FLTrust: accuracy duy trì ở 78.9\% -- minh chứng hiệu quả của cơ chế Byzantine-robust.

\subsection{So sánh Baseline (Non-FL) và Federated}
Đề tài cũng triển khai chế độ Baseline để so sánh hiệu quả liên hợp:

\begin{table}[h]
\centering
\caption{So sánh Baseline vs Federated Training (Edge-IIoT)}
\begin{tabular}{|l|l|l|}
\hline
\textbf{Chỉ số} & \textbf{Baseline (Centralized)} & \textbf{FedRL-IDS (Federated)} \\ \hline
Accuracy & 85.2\% & 83.6\% \\ \hline
F1-Macro & 82.1\% & 80.4\% \\ \hline
FPR & 0.0003 & 0.0004 \\ \hline
Privacy Leakage & Cao & Không \\ \hline
Communication Cost & -- & 6.2$\times$ reduction vs centralized \\ \hline
Byzantine Robustness & Không & 30\% malicious tolerance \\ \hline
\end{tabular}
\end{table}

Kết quả cho thấy FedRL-IDS đạt hiệu suất chỉ thấp hơn 1.6\% so với centralized baseline trong khi đảm bảo hoàn toàn quyền riêng tư dữ liệu -- mức đánh đổi chấp nhận được.

\subsection{Đánh giá RL Selector}
Phân tích hành vi của RL Selector qua 100 vòng huấn luyện cho thấy:
\begin{itemize}
    \item Early phase (Round 1--30): Selector chọn trung bình 7--8 client mỗi vòng (high exploration).
    \item Mid phase (Round 30--70): Selector học được rằng chỉ cần 4--5 client chất lượng cao là đủ.
    \item Late phase (Round 70--100): Selector chọn 3--5 client có reputation cao, tối ưu communication.
\end{itemize}
Tuy nhiên, cơ chế Diversity Penalty ngăn chặn việc chọn cùng 3 client "tốt nhất" mãi mãi (lock-in prevention), đảm bảo rằng mô hình toàn cục vẫn nhận được thông tin đa dạng từ nhiều phân phối Non-IID.

% ----------------------------------------------------------------
% 7. KẾT LUẬN VÀ HƯỚNG PHÁT TRIỂN
% ----------------------------------------------------------------
\section{Kết luận và hướng phát triển}

\subsection{Kết luận}
Dự án đã xây dựng thành công giải pháp \textbf{FedRL-IDS} toàn diện cho mạng IoT, kết hợp 5 thành phần cốt lõi hoạt động đồng bộ:
\begin{enumerate}
    \item \textbf{Kiến trúc CNN-GRU-CBAM backbone} -- trích xuất đặc trưng không gian-thời gian từ chuỗi gói tin mạng.
    \item \textbf{Thuật toán PPO} -- huấn luyện ổn định với clipped surrogate objective và focal loss.
    \item \textbf{Federated Learning} -- huấn luyện phân tán không cần tập trung dữ liệu.
    \item \textbf{FLTrust + Temporal Reputation} -- chống chịu Byzantine với ngưỡng 30\% client độc hại.
    \item \textbf{RL Client Selector} -- giảm 6.2 lần communication overhead so với centralized training.
\end{enumerate}
Việc kết hợp PPO với Federated Learning không chỉ giải quyết bài toán bảo mật quyền riêng tư mà còn mang lại độ chính xác cao (83.6\% accuracy, 80.4\% F1-Macro) trong môi trường dữ liệu Non-IID với tỷ lệ tấn công chỉ 20\%.

\subsection{Hạn chế}
\begin{itemize}
    \item Chi phí tính toán ban đầu để huấn luyện RL Selector còn cao -- cần thêm warmup rounds trước khi selector bắt đầu học hiệu quả.
    \item Mô hình nhạy cảm với việc thay đổi cấu trúc của các gói tin mạng thế hệ mới (zero-day attacks) do chỉ huấn luyện trên dữ liệu historical.
    \item Root dataset phải được cập nhật định kỳ để FLTrust duy trì hiệu quả đánh giá.
    \item Độ phức tạp tính toán tăng so với centralized training do phải huấn luyện nhiều PPO agents song song.
\end{itemize}

\subsection{Hướng phát triển}
Trong tương lai, nhóm nghiên cứu hướng tới:
\begin{enumerate}
    \item \textbf{Tích hợp AutoML:} Cơ chế tự động tối ưu hóa siêu tham số (Hyperparameter Optimization) bằng Bayesian Optimization hoặc Neural Architecture Search (NAS) để tìm kiến trúc tối ưu cho từng môi trường mạng cụ thể.
    \item \textbf{Triển khai Edge Inference:} Thực nghiệm triển khai mô hình ONNX trên các bo mạch Raspberry Pi 4 và Jetson Nano để đo lường mức tiêu thụ năng lượng, latency thực, và throughput trong điều kiện thực tế.
    \item \textbf{Dynamic Attack Adaptation:} Tích hợp cơ chế Online Learning để mô hình tự cập nhật khi phát hiện novel attack patterns mà không cần full retraining.
    \item \textbf{Đa nhiệm (Multi-Task):} Mở rộng taxonomy từ 3 lớp lên phân loại chi tiết hơn (theo attack subtype, protocol, severity) để cung cấp thông tin chi tiết hơn cho đội ngũ SOC.
    \item \textbf{Privacy Enhancement:} Tích hợp Differential Privacy ($\epsilon$-DP) vào gradient uploads để cung cấp formal privacy guarantee, thay thế cho privacy-by-design hiện tại.
\end{enumerate}



\section{TÀI NGUYÊN CỦA NHÓM}


\section{TÀI LIỆU THAM KHẢO}

\begin{thebibliography}{99}

\bibitem{fedavg}
R. McMahan, E. Moore, D. Ramage, S. Hampson, and B. A. y Arcas, ``Communication-efficient learning of deep networks from decentralized data,'' in \textit{Proc. AISTATS}, 2017.

\bibitem{rlintro}
R. S. Sutton and A. G. Barto, \textit{Reinforcement Learning: An Introduction}, 2nd ed. MIT Press, 2018.

\bibitem{fltrust}
X. Cao, M. Fang, J. Liu, and N. Z. Gong, ``FLTrust: Byzantine-robust federated learning via trust bootstrapping,'' in \textit{Proc. NDSS}, 2021.

\bibitem{fedplus}
M. El Ouadrhiri, A. Abdelhadi, M. J. Piran, and Z. Han, ``FED+: A unified approach to federated learning with personalization,'' \textit{IEEE Trans. Machine Learning in Communications and Networking}, vol. 2, pp. 580--598, 2024.

\bibitem{rludhfl}
M. Aledhari, R. Razzak, R. M. Parizi, and S. Saeed, ``RL-UDHFL: Reinforcement learning-enhanced utility-driven hierarchical federated learning for IoT,'' \textit{IEEE Internet of Things J.}, 2024.

\bibitem{fedrl_attention}
N. T. Pham, L. D. Nguyen, and N. H. Tran, ``Federated reinforcement learning based intrusion detection system using dynamic attention mechanism,'' in \textit{Proc. IEEE ICC}, 2024.

\bibitem{sfedrl}
M. A. Ferrag, O. Friha, L. Maglaras, H. Janicke, and L. Shu, ``Federated deep learning for cyber security in the internet of things: Concepts, applications, and experimental analysis,'' \textit{IEEE Access}, vol. 10, pp. 3572--3588, 2022.

\bibitem{fedrl_adaptive}
Y. Liu, J. M. Vidal, and J. Z. Liebman, ``Federated reinforcement learning for adaptive cyber defense,'' \textit{IEEE Trans. Information Forensics and Security}, 2023.

\bibitem{drl_iiot}
M. Al-Garadi, A. Mohamed, A. K. Al-Ali, X. Du, I. Ali, and M. Guizani, ``Intrusion detection approach for IIoT traffic using deep recurrent reinforcement learning assisted federated learning,'' \textit{IEEE Trans. Industrial Informatics}, vol. 19, no. 11, pp. 10964--10975, 2023.

\bibitem{prifeds}
T. D. Nguyen, P. Rieger, M. Miettinen, and A.-R. Sadeghi, ``PriFED-IDS: Privacy-preserving federated intrusion detection system,'' in \textit{Proc. ACM Wisec}, 2023.

\bibitem{drlfedrl}
N. C. Luong, D. T. Hoang, S. Gong, D. Niyato, P. Wang, Y.-C. Liang, and D. I. Kim, ``Applications of deep reinforcement learning in communications and networking: A survey,'' \textit{IEEE Communications Surveys \& Tutorials}, vol. 21, no. 4, pp. 3133--3174, 2019.

\end{thebibliography}

\end{document}